\section{Experimental Methodology}
\label{sec:method}

Each asset is split chronologically into 60\% training, 20\% validation and 20\%
sealed test on its own session index, with a one-row purge at each boundary
(Table~\ref{tab:splits}); all models see identical boundaries and dates. Hyperparameters are chosen by random search inside an
expanding-window time-series cross-validation confined to the training block:
five folds, a one-fold gap, minimum initial training fraction 0.30. Because
tuning never sees validation, validation remains a selection holdout rather than
a second tuning set. Every family gets the same budget, \TrialsUsed{} trials and at most \SecondsUsed{}~s
per cell (the protocol caps \TrialsMax{} and \SecondsMax{}~s), and every attempted
configuration is logged.

\begin{figure}[t]
  \centering
  \begin{tikzpicture}[
      node distance=5mm and 7mm,
      box/.style={draw, rounded corners=1.2pt, align=center, font=\scriptsize,
                  inner sep=3pt, minimum height=8.5mm},
      arr/.style={-{Latex[length=1.6mm]}, semithick},
      lbl/.style={font=\scriptsize\itshape, text=black!70}]
    \node[box, fill=blue!6] (raw) {raw snapshot\\{\tiny SHA-256 verified}};
    \node[box, fill=blue!6, right=of raw] (feat) {trailing features\\{\tiny 10 ratios, bars $\le t$}};
    \node[box, fill=blue!6, right=of feat] (split) {chronological split\\{\tiny 60/20/20 + purge}};
    \node[box, fill=orange!10, right=of split] (pre) {variant\\{\tiny P0 \dots\ P4}};
    \node[box, fill=orange!10, below=9mm of pre] (tune) {inner CV tuning\\{\tiny train block only}};
    \node[box, fill=green!8, left=of tune] (fit) {fit on train\\{\tiny scaler refit per fold}};
    \node[box, fill=green!8, left=of fit] (eval) {score on validation\\{\tiny vs persistence}};
    \node[box, fill=red!7, left=of eval] (test) {sealed test\\{\tiny not used}};

    \draw[arr] (raw) -- (feat);
    \draw[arr] (feat) -- (split);
    \draw[arr] (split) -- (pre);
    \draw[arr] (pre) -- (tune);
    \draw[arr] (tune) -- (fit);
    \draw[arr] (fit) -- (eval);
    \draw[arr, dashed] (eval) -- (test);
    \node[lbl, below=1mm of eval] {reported};
    \node[lbl, below=1mm of test] {one human-approved release};
  \end{tikzpicture}
  \caption{The experiment pipeline, drawn to match the implementation. The scaler
  is fitted \emph{inside} the model pipeline, so it is refitted per fold and
  cannot see validation rows. Dashed arrow: the test block is not evaluated by
  any run in this report.}
  \label{fig:pipeline}
\end{figure}

We report MAE, RMSE and $R^2$, plus \emph{skill} relative to persistence on
identical forecast dates:
\begin{equation*}
  \text{skill}_{\text{MAE}} \;=\; 1 - \frac{\text{MAE}_{\text{model}}}{\text{MAE}_{\text{naive}}}.
\end{equation*}
Skill is primary because it has a meaningful zero. $R^2$ is inflated on a trending level: our naive forecast itself attains $R^2 \approx 0.96$, so a high $R^2$
here indicates a trending series, not a good model. Absolute errors are never
compared across assets. Uncertainty uses a moving-block bootstrap over consecutive dates; overlapping blocks retain dependence, so we treat it as indicative only.
